\input{../../../stock/en-tete_v6.tex}

\newcommand{\titrechapitre}{Nombres réels}
\newcommand{\numerochapitre}{1}

\renewcommand{\headrulewidth}{0.4pt}
\renewcommand{\footrulewidth}{0.4pt}
\fancyfoot[L]{Raphaël JONTEF}
\fancyfoot[C]{\thepage}
\fancyfoot[R]{\href{https://www.alveusclub.com/}{Alveus}}
\fancyhead[L]{Mathématiques : Seconde}
\fancyhead[R]{\titrechapitre}

\begin{document}
\input{../../../stock/commands.tex}

% Page de titre custom
\setlength{\headheight}{15pt}
\begin{titlepage}
    \centering
    \vspace*{3cm}
    {\huge Chapitre \numerochapitre\par}
    {\Huge \bfseries\titrechapitre\par}
    \vspace{8cm}
    {\Large Cours rédigé par Raphaël JONTEF\par}
    \vspace{2em}
    {\small \textit{raphaeljontef@hotmail.com}\par}
    \vspace{2cm}

    {\normalsize pour un enseignement dans le cadre  \par}
    {\small d'un soutien scolaire encadré par \href{https://www.alveusclub.com/}{Alveus} \par}
    \vfill
    {\large \today\par}
\end{titlepage}

% plan du cours
\tableofcontents
\newpage


% -----------Corps du document--------------

\section{Introduction}

Un \notion{nombre réel} est un nombre qui peut être représenté par une partie entière munie d'un signe positif ou négatif, et une liste finie ou infinie de décimales \footnote{D'après \href{https://fr.wikipedia.org/wiki/Nombre_réel}{Wikipédia}.}. On note \notion{$\R$ l'ensemble des nombres réels}.

\subsection{Rappels sur la hiérarchie des ensembles de nombres}
Du plus petit au plus grand (en terme d'inclusion), on a~:
\begin{enumeratebf}
    \item \notion{$\N$ l'ensemble des entiers naturels}, c'est à dire des entiers positifs ($0,\, 1,\, 2,\, \dots$)
    \item \notion{$\Z$ l'ensemble des entiers relatifs}, c'est à dire des entiers positifs et négatifs ($\dots, -2, -1, 0, 1, 2, \dots$)
    \item \notion{$\mb{D}$ l'ensemble des nombres décimaux}, c'est à dire des nombres pouvant s'écrire avec une partie entière et une partie décimale finie (par exemple $3.14$, $-0.75$, $2.0$)
    \item \notion{$\mb{Q}$ l'ensemble des nombres rationnels}, c'est à dire des nombres pouvant s'écrire \underline{de façon irréductible} comme le quotient de deux entiers ($\frac{a}{b}$ avec $a, b \in \Z$ et bien sûr $b \neq 0$)
\end{enumeratebf}
\vspace{1em}
L'ensemble des nombres réels $\R$ contient tous les ensembles précédents. On a donc l'inclusion suivante~:
$$ \N \subset \Z \subset \mb{D} \subset \mb{Q} \subset \R $$
Cela signifie que~:
\begin{itemize}
    \item tout entier naturel est un entier relatif
    \item tout entier relatif est un nombre décimal
    \item tout nombre décimal est un nombre rationnel
    \item tout nombre rationnel est un nombre réel
\end{itemize}
\newpage
\subsection{Intérêt des nombres réels}

L'ensemble des nombres rationnels $\mb{Q}$ est déjà très vaste, mais il ne contient pas tous les nombres que l'on peut rencontrer en mathématiques. Par exemple, le nombre $\sqrt{2}$ (la racine carrée de $2$) n'est pas un nombre rationnel. En effet, il a est possible de démontrer que $\sqrt{2}$ ne peut pas s'écrire comme le quotient de deux entiers~:

\begin{demonstration}
    Si par l'absurde on suppose que $\sqrt{2}$ est rationnel, alors il existe deux entiers $a$ et $b$ tels que~:
    $$ \sqrt{2} = \frac{a}{b} $$
    Cette fraction est supposée irréductible par définition de $\mb{Q}$. En élevant au carré, on obtient~:
    $$ 2 = \frac{a^2}{b^2} $$
    Donc~:
    $$ 2b^2 = a^2 $$
    Cela signifie que $a^2$ est un multiple de $2$, donc $a$ est un multiple de $2$~: $a$ est pair. On peut ainsi écrire $a = 2k$ pour un certain entier $k$. En remplaçant dans l'équation précédente, on obtient~:
    $$ 2b^2 = (2k)^2 $$
    Donc~:
    $$ 2b^2 = 4k^2 $$
    En simplifiant par $2$, on obtient~:
    $$ b^2 = 2k^2 $$
    Ce qui signifie que $b^2$ est un multiple de $2$, donc $b$ est un multiple de $2$~: $b$ est pair.

    $a$ et $b$ sont donc tous les deux pairs, ce qui contredit le fait que la fraction $\frac{a}{b}$ soit irréductible (définition de $\mb{Q}$). On en conclut que notre hypothèse de départ est fausse, et que $\sqrt{2}$ n'est pas un nombre rationnel.
\end{demonstration}

Ainsi, pour inclure des nombres comme $\sqrt{2}$, on introduit l'ensemble des nombres réels $\R$.

\begin{remarque}{}{}
    Moralement, l'ensemble $\R$ est l'ensemble des nombres que l'on peut représenter sur une droite graduée, incluant à la fois les nombres rationnels (comme $\frac{3}{4}$, $-2$, $0.5$) et les nombres irrationnels (comme $\sqrt{2}$, $\pi$, $e$).
    C'est un fourre-tout très pratique en mathématiques !
\end{remarque}

\section{Infinis et intervalles de nombres réels}

\subsection{Infinis}
En mathématiques, on utilise souvent les symboles $+\infty$ et $-\infty$ pour représenter des concepts d'infini positif et négatif. \underline{Ces symboles ne sont pas des nombres réels}, mais entrent en jeu dans la notation des intervalles.

\subsection{Intervalles de nombres réels} 
Un \notion{intervalle de $\R$} est un ensemble de nombres réels compris entre deux bornes.
\begin{definition}{}{intervalle}
    On dit qu'un ensemble $I \subset \R$ est un \notion{intervalle} de $\R$ si pour tous réels $a<b$ dans $I$, tout réel $x$ compris strictement entre $a$ et $b$ appartient aussi à $I$. \\\\Moralement, $I$ n'est pas troué : si deux de ses éléments sont choisis, tous les réels entre ces deux éléments font aussi partie de $I$.
\end{definition}

\begin{exemple}{}{d'intervalles de $\R$ et leur notation}
    Les ensembles suivants sont des intervalles de $\R$~:
    \begin{itemize}
        \item $[2, 5] = \{ x \in \R \mid 2 \leq x \leq 5 \}$ (intervalle fermé)
        \item $] -3, 1 [ = \{ x \in \R \mid -3 < x < 1 \}$ (intervalle ouvert)
        \item $[0, +\infty[ = \{ x \in \R \mid x \geq 0 \}$ (intervalle semi-ouvert)
        \item $]-\infty, 4] = \{ x \in \R \mid x \leq 4 \}$ (intervalle semi-ouvert)
        \item $\R = ]-\infty, +\infty[$ (l'unique intervalle de $\R$ non borné : lui-même !)
    \end{itemize}
\end{exemple}

\begin{remarque}{}{concernant les bornes}
    Dans la notation des intervalles, les crochets fermants $[$ et $]$ signifient que la borne est \underline{incluse} dans l'intervalle (on dit que l'intervalle est \notion{fermé} en cette borne). \\
    À l'inverse, les crochets $]$ et $[$ signifient que la borne est \underline{exclue} de l'intervalle (on dit que l'intervalle est \notion{ouvert} en cette borne).\\\\
    On devine ainsi pourquoi $+\infty$ et $-\infty$ sont toujours accompagnés de crochets ouverts dans la notation des intervalles : on ne peut pas inclure l'infini dans un intervalle, puisqu'il ne s'agit pas d'un nombre réel !
\end{remarque}

\begin{exemple}{}{à faire soi-même}
    Caractériser les notations des intervalles suivants, à la façon de l'ensemble précédent~:
    \begin{itemize} 
        \item $[-1, 3[$
        \item $]2, +\infty[$
        \item $]-\infty, 0]$
    \end{itemize}
\end{exemple}

\begin{exemple}{}{à faire soi-même}
    Dans l'autre sens, simplifier les notations de ces ensembles~:
    \begin{itemize} 
        \item $\{ x \in \R \mid -5 \leq x \leq 2 \}$
        \item $\{ x \in \R \mid x > 7 \}$
        \item $\{ x \in \R \mid x \leq -1 \text{ ou } x \geq 3 \}$
    \end{itemize}
\end{exemple}

\end{document}